\section{The complete record: endpoint work and radius--roof disintegration}
\label{sec:rc-complete}
The preceding clock calculation extends beyond reversing records, but
its endpoint work must then be retained.  This section treats every
regular physical $m$-flight record, with all original exact labels,
not just the critical tubes of Section~\ref{sec:pt-roof}.  It produces
a transverse coordinate for a parameter-extended source.  It does
not change the fixed-radius probability space or the raw target.

\begin{theorem}[Uniform continuation at prescribed endpoint momenta]
\label{thm:rc-continuation}
Fix any actual regular record and its two endpoint momenta
$(a,v)=(p_0,p_m)\in(-1,1)^2$.  Keep them fixed while varying the
common radius.  Replace the endpoint constants in
\eqref{eq:cc-reflection} by
\[
 E_0=-\arcsin a\pmod{2\pi},\qquad
 E_m=\pi+\arcsin v\pmod{2\pi}.
\]
The full derivative matrix is still \eqref{eq:cc-matrix}, now with
$c_0=\sqrt{1-a^2}$ and $c_m=\sqrt{1-v^2}$ in place of one.
All three bounds \eqref{eq:cc-inverse} remain valid, as do the common
formal radius disk \eqref{eq:cc-analytic-constants} and the physical
radius response \eqref{eq:cc-radius-derivative}.  None of their
constants depends on $(a,v)$, the word, or the smallest incidence.
Physicality is asserted only on connected regular physical
continuation intervals, not on the whole formal radius disk.
\end{theorem}
\begin{proof}
The departure direction is $\alpha_0+\arcsin a$.  The incoming
terminal direction is $\alpha_m+\pi-\arcsin v$; tangential momentum
is unchanged by the last specular reflection.  These give exactly
the two displayed constants modulo $2\pi$.  Their derivatives with
respect to $R$ vanish because $a,v$ are fixed.
The flight differentiation in Theorem~\ref{thm:cc-full-resolvent}
then gives the same matrix with its true endpoint cosines.  Its
column excess remains exactly $b_i$, including the endpoint columns:
each endpoint cosine appears both in its diagonal correction and in
the sole off-diagonal entry of that column.  The unweighted and
weighted column arguments therefore apply without alteration, and
the row bound follows by the same geometric sum.  No derivative of
$\arcsin a$ or $\arcsin v$ is taken.  The complex flight estimates
in Theorem~\ref{thm:cc-radius} use only the positive flight lengths,
the disk radius bound and this inverse, not normality.  They supply
the same formal continuation and the same bound on $\alpha'$.
The interpretation of its physical subinterval is unchanged.
\end{proof}

\begin{theorem}[The physical roof has only bounded endpoint work]
\label{thm:rc-clock}
Let $L(R)$ be the true physical roof on any connected physical
continuation interval from Theorem~\ref{thm:rc-continuation}.
Then
\begin{equation}\label{eq:rc-exact-work}
 L'(R)=-c_0-c_m-2\sum_{i=1}^{m-1}c_i
             +R\{-a\alpha_0'(R)+v\alpha_m'(R)\}.
\end{equation}
Define the collision-independent constants
\[
 C_{\rm end}=2R_+D_R,\qquad C_{\rm loss}=C_{\rm end}+2\gamma_0,
 \qquad m_0=\left\lceil2C_{\rm loss}/\kappa_0\right\rceil.
\]
For every $m\geq1$ and every such record,
\begin{equation}\label{eq:rc-full-slope}
 \kappa_0m-C_{\rm loss}\leq -L'(R)\leq2m+C_{\rm end}.
\end{equation}
In particular for all $m\geq m_0$, the actual roof is strictly
monotone on each connected physical component, and
$-L'(R)\geq\kappa_0m/2$.  These statements are uniform over the
entire original radius interval, over the endpoint momenta and over
all physical collision words.  They impose no new incidence or
clearance guard on the original source.  The conservative constant
$m_0$ is a theoretical threshold, not an experimentally reached
collision count.
\end{theorem}
\begin{proof}
On the stationary interior of the full action, its endpoint gradient
is $R(-a,v)$.  Differentiating the action along the prescribed-momentum
continuation consequently leaves the two endpoint chain-rule terms,
in addition to the explicit radius derivative.  The latter telescopes
as in \eqref{eq:cc-hadamard}, proving \eqref{eq:rc-exact-work}.
Its endpoint work has absolute value at most
$R(|a|+|v|)D_R\leq C_{\rm end}$.
The visited incidence sum is bounded below by
$\gamma_0\lfloor m/94\rfloor$ from
Corollary~\ref{cor:cb-collision-bending}.  Dropping only the two
nonnegative endpoint cosines in a lower bound, and using
$\lfloor m/94\rfloor\geq m/94-1$, gives
\[
 -L'(R)\geq 2\gamma_0\lfloor m/94\rfloor-C_{\rm end}
          \geq\kappa_0m-C_{\rm loss}.
\]
There are exactly $2m$ incidence contributions with multiplicity,
each at most one, which gives the upper bound.  The definition of
$m_0$ proves the last assertion.  This is a pointwise geometric
estimate for every physical record, not an ergodic average.  In
particular no independence, pressure input or collision-word count
occurs in its proof.
\end{proof}

\begin{theorem}[An exact disintegration of the parameter-extended physical source]
\label{thm:rc-disintegration}
Fix $m\geq m_0$ and a nonnegative Borel function $h(R,x)$ on the
regular physical collision space, zero off the original initial
section when section probability is used.  It may include the exact
original selector $\mathbf1_{\Pi_{n,k,m,R}}$, with the true section
visits evaluated at that radius.  For a regular root with endpoint
momenta $(a,v)$ let
\[
 u(R,x)=\left|\partial_{\alpha_0}p_m\right|^{-1}
        \quad\text{with $R,p_0=a$ fixed}.
\]
Use the physical momentum branches and one-copy convention of
Section~\ref{sec:mc-fibres}.  For every nonnegative Borel test $\phi$,
\begin{equation}\label{eq:rc-disintegration}
 \begin{split}
 &\int_{R_-}^{R_+}\int
       h(R,x)\phi(L_{m,R}(x))\,d\nu_R^*(x)\,dR\\
 &=\frac1{4\pi c_*}\int_{-1}^1\int_{-1}^1\int_{\R}
       \phi(t)
       \sum_{\substack{(R,x):\ p_0=a,\ p_m=v\\L_{m,R}(x)=t}}
       \frac{h(R,x)u(R,x)}{|L'_{a,v,x}(R)|}\,dt\,dv\,da .
 \end{split}
\end{equation}
The derivative is the full derivative with the endpoint work in
\eqref{eq:rc-exact-work}, not a derivative at fixed initial angle.
The sums include only genuine regular physical records and count
each once.  No finiteness of either side is assumed.  The new
Jacobian factor satisfies
\begin{equation}\label{eq:rc-jacobian-bounds}
 \frac1{2m+C_{\rm end}}
  \leq\frac1{|L'_{a,v,x}(R)|}\leq\frac2{\kappa_0m}.
\end{equation}
The equality is auxiliary geometric information in the added radius
coordinate.  It makes no assertion that the radius integral can
replace the required supremum over $R$.
\end{theorem}
\begin{proof}
On a regular physical branch the initial collision measure is
$d\alpha_0\,da/(4\pi)$ and the map to $(a,v)$ has positive absolute
Jacobian $|\partial_{\alpha_0}p_m|$; this is the actual root
Jacobian used in the physical momentum change of variables in
Section~\ref{sec:mc-fibres}.  After the original section restriction
its density is $u/(4\pi c_*)$.  This is a local change of variables;
the all-time upper kernel bound in that section is not used here.
In particular it is not a composition of one-step momentum kernels.

Theorem~\ref{thm:rc-continuation} supplies local physical branches
parametrized by $(R,a,v)$.  For $m\geq m_0$ the coordinate change
$(R,a,v)\mapsto(L(R,a,v),a,v)$ is locally invertible, with absolute
determinant $|L'_{a,v,x}(R)|$, by Theorem~\ref{thm:rc-clock}.
The ordinary substitution theorem therefore gives
\eqref{eq:rc-disintegration} on one chart, with the reciprocal
Jacobian in the stated direction.  Countably many physical charts
cover the regular source: use countably many lattice words and
second countability in $(R,\alpha_0,a)$.  A fixed first-chart
allocation makes their representations disjoint.  The allocation
indicators and any section or exact-label cuts stay inside $h$;
none is differentiated.  Applying substitution to finite positive
chart sums and then monotone convergence gives the equality for
the complete regular source, allowing infinite values.  The
identity concerns the regular source and makes no claim about the
separate paired-boundary traces.
Finally \eqref{eq:rc-full-slope} proves
\eqref{eq:rc-jacobian-bounds}.
\end{proof}

\begin{remark}[The precise unpaid step is unchanged]
\label{rem:rc-unpaid}
This exact formula retains all regular physical records, all
original collision labels and their true physical roofs.  It is
not a new probability model and does not discard low-incidence
records.  Nevertheless its crossing sum is taken at different
radii.  The unscaled root-kernel bound at a common $(R,a,v)$ does
not control this sum.  Nor does an estimate integrated in $R$
control the original fixed-radius roof essential height.  A
uniform bound for those concentrations, the complete positive
source and paired traces, or the original comparison capacity
is still required to prove \eqref{eq:v112-remaining}.  This remains
a substantive mathematical obligation rather than an automatic
consequence of the factor $1/m$ in a different-coordinate
Jacobian.  All original probability quantifiers and ordered
microscopic limits are retained.
\end{remark}
